\section{What we built}
\label{sec:system}

The task is news recommendation from click-logs. The unit of everything is an \term{impression}:
at time $t$ a user is shown a list of \term{candidate} articles and clicks some of them. Our job is
to order each impression's candidates so the clicked ones come first. We extend our two
Assignment-1 retrieval systems (lexical and semantic) with the \term{behavioural} signal in the
logs: what users clicked, how often each article has been clicked or shown, and how fresh it is.

\subsection{Data and evaluation protocol}

\begin{table*}[t]
\centering\footnotesize
\begin{tabular}{@{}lrcrr@{}}
\toprule
 & & & \multicolumn{2}{c}{leaderboard test set} \\
\cmidrule(l){4-5}
dataset & articles & train / val / test impressions & impressions & articles \\
\midrule
\eb{} small (Danish) & 20,738 & 168,522 / 64,365 / 244,647 & 13,536,710 & 125,500 \\
\mind{} small (English) & 65,238 & 95,071 / 61,894 / 73,152 & 2,370,727 & 120,961 \\
\bottomrule
\end{tabular}
\caption{Development data and the held-out leaderboard sets. Splits are temporal: every train
impression precedes every val impression, which precedes every test impression.}
\label{tab:data}
\end{table*}

\textbf{Splits are by time, never random.} A random split would let the model learn from
Thursday to predict Wednesday, which no deployed system can do. \eb{} val begins 23 May 07:00 and
test 25 May 07:00; \mind{} uses the last day of its training log as val and its dev log as test.

\textbf{Metric: per-impression AUC}, computed inside each impression and averaged over those with
both a click and a non-click. We never pool candidates across impressions, because the ranking is
only judged inside one: pooling rewards features that merely differ between impressions, and the
effect is large, with article freshness reading 0.5861 pooled against 0.5087 per impression.

\textbf{Significance: paired bootstrap 95\% intervals} \cite{efron1993}, 1,000 resamples, both
systems scored on the \emph{same} resampled impressions so their correlated errors cancel. A gain
counts only if its interval excludes zero. Selection happens on val; test is scored once.

\subsection{Architecture: retrieve, then re-rank}

\begin{figure*}[t]
\centering
\begin{tikzpicture}[
  box/.style={draw, rounded corners=2pt, align=center, font=\footnotesize, inner sep=2.5pt,
              minimum height=7mm, text width=#1},
  arr/.style={-{Stealth[length=1.8mm]}, semithick}]
\node[box=2.2cm, fill=gray!8] (imp) at (0,0) {impression at time $t$: history, candidates};
\node[box=2.2cm, fill=blue!6] (bm25) at (3.2,0.62) {BM25 over history titles};
\node[box=2.2cm, fill=blue!6] (ann) at (3.2,-0.62) {exact kNN on mean history vector};
\node[box=1.5cm, fill=blue!6] (fuse) at (5.75,0) {fuse, keep top 200};
\node[box=2.3cm, fill=orange!10] (feat) at (8.3,0) {22 causal features (data before $t$)};
\node[box=1.7cm, fill=orange!10] (lgb) at (10.8,0) {LightGBM LambdaRank};
\node[box=1.15cm, fill=gray!8] (out) at (12.75,0) {ranked list};
\draw[arr] (imp.east) -- ++(0.2,0) |- (bm25.west);
\draw[arr] (imp.east) -- ++(0.2,0) |- (ann.west);
\draw[arr] (bm25.east) -- ++(0.2,0) |- (fuse.west);
\draw[arr] (ann.east) -- ++(0.2,0) |- (fuse.west);
\draw[arr] (fuse) -- (feat);
\draw[arr] (feat) -- (lgb);
\draw[arr] (lgb) -- (out);
\node[draw=blue!45, dashed, rounded corners=3pt, fit=(bm25)(ann)(fuse), inner sep=3pt,
      label={[font=\scriptsize]below:stage one: cheap, whole catalogue}] {};
\node[draw=orange!70, dashed, rounded corners=3pt, fit=(feat)(lgb), inner sep=3pt,
      label={[font=\scriptsize]below:stage two: rich, candidates only}] {};
\end{tikzpicture}
\caption{The two-stage pipeline. Stage one touches every article, so it must stay cheap per
article. Stage two touches only the candidates, so it can afford per-candidate features.}
\label{fig:pipeline}
\end{figure*}

\textbf{Stage one, candidate generation} (carried from Assignment 1). \emph{BM25}
\cite{robertson2009,lu2024bm25s} is a lexical score: the query is the titles of the user's last
$N$ clicks ($N=30$ on \eb{}, $100$ on \mind{}), and an article scores highly when it shares rare
words with them. Stemming is on for both languages and a Danish stopword list is applied. The
\emph{semantic} retriever represents the user as the mean of their clicked articles' vectors and
finds the nearest articles by cosine with an exact FAISS \texttt{IndexFlatIP} \cite{johnson2019faiss}.
Vectors are the provided 768-dimensional \texttt{contrastive\_vector} for \eb{} and
384-dimensional \texttt{all-MiniLM-L6-v2} \cite{reimers2019} for \mind{}, both chosen by measurement
in Assignment 1. The two scores are standardised within each candidate pool and mixed linearly,
with the weight chosen on val. Each retriever keeps its top 200 candidates. Both indexes are held
in RAM.

\textbf{Stage two, the re-ranker.} A LightGBM \cite{ke2017lightgbm} gradient-boosted tree model
with the \term{LambdaRank} objective \cite{burges2010}. LambdaRank learns from pairs of candidates
inside the same impression and weights each pair by how much swapping them would change nDCG, so
training optimises the ranking itself rather than a click probability. We measured this choice
rather than assuming it (Table~\ref{tab:choices}). Settings: learning rate 0.05, 31 leaves,
at least 50 rows per leaf, 90\% feature and row subsampling, up to 400 rounds with early stopping
on val nDCG, seed 13.

\subsection{Features and the behaviour-window boundary}

\textbf{Every feature for an impression at time $t$ uses only data from strictly before $t$.} We
call this the \term{behaviour-window boundary}. Click and exposure counts are computed by binary
search over each article's sorted event times with a strict ``$< t$'' cut, so an event at exactly
$t$ is excluded. Tests assert the boundary on every split, and each test was checked by injecting a
future click and confirming it fails (Section~\ref{sec:q9}).

\begin{table*}[t]
\centering\small
\begin{tabularx}{\linewidth}{@{}l>{\raggedright\arraybackslash}Xcc@{}}
\toprule
family & features and what they measure & \eb{} & \mind{} \\
\midrule
stage-one scores & \feat{bm25}, \feat{emb} (the two retrieval scores); \feat{emb_rank} (the
embedding score's rank inside the impression) & yes & yes \\
click popularity & \feat{pop_causal}: clicks before $t$ in a 6\,h window (unbounded on \mind{});
\feat{pop_24h}: the same over 24\,h; \feat{pop_velocity}: their ratio, ``rising now'' versus
``steadily popular''; \feat{pop_rank}, \feat{pop_rel_max}: standing among this impression's
candidates & yes & yes \\
exposure & the same five shapes (\feat{exp_causal}, \feat{exp_24h}, \feat{exp_velocity},
\feat{exp_rank}, \feat{exp_rel_max}) counting how often an article was \emph{shown} before $t$ &
yes & yes \\
user match & \feat{cat_match}: share of the user's history in the candidate's category;
\feat{ent_overlap}: shared named entities with the history & yes & yes \\
context & \feat{hist_len}: history length; \feat{n_cands}: number of candidates in the impression &
yes & yes \\
freshness & \feat{age_hours}, \feat{recency} (24\,h half-life decay of age) & yes & no \\
engagement & \feat{user_read}, \feat{user_scroll}: mean read time and scroll depth of past clicks;
\feat{engage_sim}: cosine to a read-time-weighted user vector & yes & no \\
\midrule
total & & 22 & 17 \\
\bottomrule
\end{tabularx}
\caption{The re-ranker's features. Rank and relative features exist because the metric is decided
inside one impression: a popularity of 200 dominates a quiet impression and is ordinary in a busy
one.}
\label{tab:features}
\end{table*}

\textbf{Five features cannot exist on \mind{}, and we drop them explicitly.} \mind{}'s publication
time is null for all 65,238 articles, and it ships no read-time or scroll data at all. A column of
zeros looks like a real feature to LightGBM, so the builder refuses to emit any feature that is
all-null or constant. This matches the data audit we did before designing features: \eb{} is rich
in behavioural signal, \mind{} has little beyond the click itself.

\textbf{Exposure} (``how often was this article shown before $t$'') was added for a practical
reason and kept for a better one. The leaderboard test files withhold the clicks, so click
popularity cannot be computed there (Section~\ref{sec:leaderboard}), but the lists of shown
articles are present. A live recommender also knows its own exposure log immediately, while click
feedback arrives later, so exposure is arguably the more realistic signal.
